%-------------------------
% Resume in Latex
% Author : Stéphane KPOVIESSI
% Target: Big Data & AI Engineer
% Format: Two-Page Resume
%-------------------------

\documentclass[letterpaper,11pt]{article}

\usepackage{latexsym}
\usepackage[empty]{fullpage}
\usepackage{titlesec}
\usepackage[usenames,dvipsnames]{color}
\usepackage{enumitem}
\usepackage[colorlinks=true, urlcolor=primary, linkcolor=primary]{hyperref}
\usepackage{fancyhdr}
\usepackage[french]{babel}
\usepackage{tabularx}
\usepackage{needspace}
\usepackage{graphicx}
\usepackage{tikz}
\input{glyphtounicode}

% --- COULEURS ---
\definecolor{primary}{RGB}{0, 79, 144}
\definecolor{darkgrey}{RGB}{50, 50, 50}

%----------FONT OPTIONS----------
\usepackage[sfdefault]{noto-sans}
\usepackage[T1]{fontenc}

\pagestyle{fancy}
\fancyhf{}
\fancyfoot{}
\renewcommand{\headrulewidth}{0pt}
\renewcommand{\footrulewidth}{0pt}
\setlength{\footskip}{4.1pt}

% --- MARGES ---
\addtolength{\oddsidemargin}{-0.5in}
\addtolength{\evensidemargin}{-0.5in}
\addtolength{\textwidth}{1.0in}
\addtolength{\topmargin}{-0.65in}
\addtolength{\textheight}{1.45in}

\urlstyle{same}
\raggedbottom
\raggedright
\setlength{\tabcolsep}{0in}

% --- FORMATAGE ---
\titleformat{\section}{
  \vspace{-4pt}\scshape\raggedright\large\color{primary}
}{}{0em}{}[\color{primary}\titlerule \vspace{0pt}]

\pdfgentounicode=1

%-------------------------
% Custom commands
\newcommand{\resumeItem}[1]{
  \item\small{{#1 \vspace{-1.5pt}}}
}

\newcommand{\resumeSubheading}[4]{
  \vspace{1pt}\item
    \begin{tabular*}{0.97\textwidth}[t]{l@{\extracolsep{\fill}}r}
      {\color{primary}\textbf{#1}} & #2 \\
      \textit{\small#3} & \textit{\small #4} \\
    \end{tabular*}\vspace{-4pt}
}

\newcommand{\resumeProjectHeading}[2]{
    \Needspace{6\baselineskip}
    \vspace{2pt}\item
    \begin{tabular*}{0.97\textwidth}{l@{\extracolsep{\fill}}r}
      \small#1 & #2 \\
    \end{tabular*}\vspace{-3pt}
}

\newcommand{\resumeSubItem}[1]{\resumeItem{#1}\vspace{-4pt}}
\renewcommand\labelitemii{$\vcenter{\hbox{\tiny$\bullet$}}$}
\newcommand{\resumeSubHeadingListStart}{\begin{itemize}[leftmargin=0.15in, label={}]}
\newcommand{\resumeSubHeadingListEnd}{\end{itemize}}
\newcommand{\resumeItemListStart}{\begin{itemize}[leftmargin=0.15in, label={\small\textbullet}, itemsep=1.5pt, parsep=0pt, topsep=2pt, partopsep=0pt]}
\newcommand{\resumeItemListEnd}{\end{itemize}\vspace{-1pt}}

%-------------------------------------------
\begin{document}

%----------HEADING----------
\noindent
\begin{minipage}[c]{0.20\textwidth}
  \centering
  \begin{tikzpicture}
    \begin{scope}
      \clip (0,0) circle (1.45cm);
      \node at (0,0) {\includegraphics[width=2.9cm]{./profile.jpeg}};
    \end{scope}
    \draw[color=primary, line width=1.2pt] (0,0) circle (1.45cm);
  \end{tikzpicture}
\end{minipage}%
\hspace{0.01\textwidth}%
\begin{minipage}[c]{0.77\textwidth}
  \centering
  {\Huge \scshape \color{primary} Stéphane KPOVIESSI} \\[3pt]
  {\large \textbf{Ingénieur Big Data \& Intelligence Artificielle}} \\[7pt]
  \small
  Île-de-France (Cergy)
    \enspace$|$\enspace Disponible en France et à l'international
    \enspace$|$\enspace +33\,7\,48\,59\,47\,13 \\[2pt]
  \href{mailto:oastephanekpoviessi@gmail.com}{oastephanekpoviessi@gmail.com}
    \enspace$|$\enspace
  \href{https://linkedin.com/in/stephanekpoviessi}{LinkedIn}
    \enspace$|$\enspace
  \href{https://github.com/EngineerProjects}{GitHub}
    \enspace$|$\enspace
  \href{https://seshat-ai.com}{Seshat-AI}
    \enspace$|$\enspace
  \href{https://kpoviessi-stephane.vercel.app/}{Portfolio}
\end{minipage}
\vspace{3pt}

%-----------PROFIL-----------
\section{Profil}
\small{
Ingénieur Big Data \& Data Science spécialisé en IA générative, systèmes multi-agents et industrialisation Data/IA. Expérience en environnement critique chez Allianz sur Python/PySpark et Azure Databricks, complétée par la conception de systèmes GenAI avancés intégrant RAG hybride, orchestration multi-agents, LLM multi-fournisseurs et automatisation de workflows métiers. Bonnes pratiques d'ingénierie logicielle appliquées en Python/Go, APIs, Docker et CI/CD.}
\vspace{4pt}


%-----------EXPERIENCES-----------
\section{Expériences Professionnelles}
\resumeSubHeadingListStart

\resumeSubheading
{Ingénieur Data / Data Scientist}{Nov. 2025 -- Avr. 2026}
{Allianz France --- Direction P4D, Équipe Advanced Data \& Climate}{Paris, France}
\resumeItemListStart
  \resumeItem{Migration d'un datamart Assurance de SAS vers \textbf{Python/PySpark}, avec garantie de la parité fonctionnelle sur plus de 15 000 lignes de code.}
  \resumeItem{Implémentation d'une architecture \textbf{Médaillon} sur Azure Databricks intégrant plus de 150 règles de gestion métier pour fiabiliser l'alimentation de cas d'usage analytiques.}
  \resumeItem{Conception et automatisation d'un pipeline de \textbf{contrôle qualité des données} géospatiales pour renforcer la robustesse du reporting et des traitements aval.}
  \resumeItem{Développement d'un framework de \textbf{tests automatiques} SAS/Python pour sécuriser la non-régression, l'industrialisation et la mise en production.}
\resumeItemListEnd

\resumeSubheading
{Data Analyst -- Missions d'appui stratégique}{2023 -- 2025}
{LYNE ET FRERES SARL \normalfont\small(Missions récurrentes d'analyse Data)}{Ikpinlè, Bénin / À distance}
\resumeItemListStart
  \resumeItem{Analyse de données commerciales, opérationnelles et de marché avec \textbf{Python/Pandas} pour soutenir les décisions de diversification et le suivi des indicateurs clés de performance.}
  \resumeItem{Conception de rapports et tableaux de bord interactifs avec \textbf{Streamlit} et interfaces web HTML/CSS/JavaScript pour restituer les analyses au personnel de l'entreprise.}
  \resumeItem{Suivi analytique des activités agricoles, notamment les plantations de palmier et les étapes de transformation, afin d'identifier des pistes d'amélioration opérationnelle.}
  \resumeItem{Mise à disposition et accompagnement à l'usage responsable d'outils IA internes basés sur \textbf{OpenWebUI} et modèles open source pour appuyer les usages métiers.}
\resumeItemListEnd

\resumeSubheading
{Ingénieur Machine Learning}{Mai 2024 -- Août 2024}
{Groupe Sylvagreg \normalfont\small(Stage)}{Lille, France} 
\resumeItemListStart
  \resumeItem{Recherche et développement d'une solution de reconstruction 3D à partir d'images 2D, combinant Computer Vision, photogrammétrie et modèles génératifs sous PyTorch.}
  \resumeItem{Automatisation des pipelines de prétraitement, d'augmentation de données et de génération pour fiabiliser les traitements d'inférence.}
  \resumeItem{Déploiement cloud sur \textbf{AWS} avec conteneurisation \textbf{Docker} et interface de visualisation pour les utilisateurs métiers.}
\resumeItemListEnd

\resumeSubHeadingListEnd

%-----------PROJECTS-----------
\section{Projets Sélectionnés}
\resumeSubHeadingListStart

\resumeProjectHeading
{\textbf{Seshat \& Seshat-AI} $|$ \emph{Go, TypeScript, Electron, PostgreSQL, RAG, MCP, Agents IA}}{\href{https://github.com/EngineerProjects/seshat}{GitHub}}
\resumeItemListStart
  \resumeItem{Conception et développement de bout en bout d'une plateforme d'agents IA multi-dépôts : moteur Go, backend local embarqué, backend cloud multi-tenant, application desktop Electron/React et console d'administration.}
  \resumeItem{Conception d'un moteur \textbf{RAG hybride} (chunking, embeddings multi-fournisseurs, recherche vectorielle) et d'un orchestrateur \textbf{multi-agents} exécutant des outils externes via le protocole \textbf{MCP}.}
  \resumeItem{Architecture d'une couche d'abstraction multi-fournisseurs LLM (OpenAI, Mistral, Gemini, Ollama) avec héritage automatique des capacités entre services (texte, embeddings, image, audio).}
  \resumeItem{Développement d'une application desktop sécurisée : identifiants chiffrés via le trousseau système (DPAPI/Keychain/libsecret), authentification et gestion de sessions.}
  \resumeItem{Conception d'une architecture cloud multi-tenant orientée entreprise (permissions, audit, supervision human-in-the-loop, tâches planifiées) avec automatisations métier réelles : synchronisation Gmail/WhatsApp, prospection commerciale.}
\resumeItemListEnd

\resumeProjectHeading
{\textbf{Tech Watch Agent} $|$ \emph{Python, FastAPI, LangGraph, PostgreSQL, pgvector, React, Docker}}{\href{https://github.com/EngineerProjects/tech-watch-agent}{GitHub}}
\resumeItemListStart
  \resumeItem{Développement d'une plateforme de \textbf{veille technologique multi-agents} automatisant la planification, la collecte multi-source et la synthèse de rapports ciblés.}
  \resumeItem{Backend FastAPI avec orchestration LangGraph et mémoire sémantique PostgreSQL/\textbf{pgvector}, industrialisé en stack Docker avec frontend React.}
\resumeItemListEnd

\resumeProjectHeading
{\textbf{BI Retail Solution} $|$ \emph{Python, PySpark, Airflow, PostgreSQL, Power BI, Docker}}
{\href{https://github.com/EngineerProjects/BI_Retail}{GitHub}}
\resumeItemListStart
    \resumeItem{Mise en place d'une chaîne Data pour l'exploitation de données de ventes et de supply chain : génération et ingestion des données, transformations \textbf{PySpark} et chargement dans un \textbf{Data Warehouse PostgreSQL}.}
    \resumeItem{Orchestration des traitements avec \textbf{Apache Airflow} et déploiement de l'environnement PostgreSQL/Airflow avec \textbf{Docker Compose}.}
    \resumeItem{Conception de tableaux de bord \textbf{Power BI} pour l'analyse des ventes, des clients et de la performance logistique : revenus, stocks, coûts d'expédition, profits et performance par produit, pays et magasin.}
\resumeItemListEnd

\resumeProjectHeading
{\textbf{Solar Irradiance Forecasting} $|$ \emph{Python, Pandas, Scikit-learn, FastAPI, REST}}
{\href{https://github.com/EngineerProjects/solar-forecasting}{GitHub}}
\resumeItemListStart
  \resumeItem{Conception d'un pipeline de prévision de l'irradiance solaire à 24h : ingestion de données météorologiques via l'API Open-Meteo, validation de la qualité des données et feature engineering temporel.}
  \resumeItem{Entraînement et évaluation d'un modèle \textbf{Random Forest} ($R^2$ = 0,85), avec tests unitaires et exposition des prédictions via une \textbf{API REST FastAPI}.}
\resumeItemListEnd

\resumeSubHeadingListEnd


%-----------SKILLS-----------
\section{Compétences}
\begin{itemize}[leftmargin=0.15in, label={}]
\small{\item{
\textbf{Programmation} : Python, SQL, Go, R, Bash, C \\
\textbf{Automation \& IA} : LLM, RAG, NLP, chatbots, systèmes multi-agents, LangGraph, LangChain, automatisation de workflows, human-in-the-loop \\
\textbf{Data Engineering \& Analytics} : PySpark, Delta Lake, ETL/ELT, dbt, Apache Airflow, Azure Databricks, architecture Médaillon, Data Quality, Power BI \\
\textbf{Machine Learning \& Deep Learning} : Pandas, Scikit-learn, XGBoost, PyTorch, TensorFlow, Keras, CNN, Vision Transformers, forecasting, évaluation de modèles \\
\textbf{Backend \& APIs} : FastAPI, REST APIs, SSE, PostgreSQL, pgvector, intégration d'outils externes \\
\textbf{Outils \& Cloud} : Git, Docker, CI/CD, monitoring, Azure (ADLS, ADF, Databricks), AWS (EC2, S3), Streamlit, React
}}
\end{itemize}



%-----------FORMATION-----------
\section{Formation}
\resumeSubHeadingListStart
  \resumeSubheading
    {JUNIA ISEN}{Lille, France}
    {Diplôme d'Ingénieur --- Big Data \& Data Science}{Diplômé en 2026}
  \resumeSubheading
    {JUNIA ISEN}{Rabat, Maroc}
    {Licence --- Électronique et Systèmes Numériques}{2023}
  \resumeSubheading
    {Cycle préparatoire Marie Stella}{Cotonou, Bénin}
    {Classe préparatoire MPSI --- Mathématiques, Physique, Sciences de l'Ingénieur}{2022}
\resumeSubHeadingListEnd

%-----------LANGUES-----------
\section{Langues}
\begin{itemize}[leftmargin=0.15in, label={}]
  \small{\item{
   \textbf{Français} : Langue maternelle \vspace{1.5pt} \\
   \textbf{Anglais} : Niveau professionnel (B2 - Certifié Linguaskill Business English) \vspace{1.5pt} \\
   \textbf{Yoruba, Fon, Goun} : Langues maternelles (Bilingue)
  }}
\end{itemize}

%-----------ENGAGEMENTS-----------
\section{Engagements \& Associations}
\resumeSubHeadingListStart
  \resumeSubheading
    {Vice-Président}{Sep. 2024 -- Août 2025}
    {Commission Réseau JUNIA ALUMNI}{Lille, France}
    \resumeItemListStart
      \resumeItem{Coordination de la commission, animation du réseau d'alumni et contribution à l'organisation d'afterworks et d'événements de networking.}
      \resumeItem{Représentation de l'association auprès des étudiants et coordination avec le Bureau des Étudiants (BDE) pour fluidifier les actions communes.}
    \resumeItemListEnd

  \resumeSubheading
    {Responsable Communication}{Sep. 2023 -- Août 2024}
    {Commission Réseau JUNIA ALUMNI}{Lille, France}
    \resumeItemListStart
      \resumeItem{Pilotage de la communication de l'association : préparation des annonces, valorisation des événements et diffusion des informations auprès de la communauté.}
      \resumeItem{Création de supports de communication et animation de la visibilité des initiatives pour renforcer l'engagement du réseau alumni.}
    \resumeItemListEnd
\resumeSubHeadingListEnd

\end{document}
